\documentclass[11pt,a4paper,spanish]{article}
\usepackage[spanish]{babel}
\usepackage{fontspec}
\usepackage[left=2cm,right=2cm,top=2cm,bottom=2cm]{geometry}
\usepackage{amsmath}
\usepackage{amssymb}
\usepackage{graphicx}
\usepackage{array}
\usepackage{booktabs}
\usepackage{hyperref}
\usepackage{float}

\title{\textbf{XGBoost para Análisis de Riesgo Crediticio} \\ {\large Guía Explicativa Paso a Paso}}
\author{Material Educativo - Departamento de Riesgo}
\date{\today}

\begin{document}

\maketitle

\begin{abstract}
Este documento explica de manera clara y accesible cómo funciona el algoritmo XGBoost (Extreme Gradient Boosting) en el contexto del análisis de riesgo crediticio. Se incluyen ejemplos prácticos con datos reales, explicaciones teóricas simplificadas y flujos de datos completos para que cualquier persona pueda entender cómo el modelo toma decisiones sobre quién es un riesgo de impago.
\end{abstract}

\tableofcontents
\newpage

% ============================================================================
\section{Introducción: ¿Qué es XGBoost?}
% ============================================================================

XGBoost es un algoritmo de \textbf{aprendizaje automático} diseñado para hacer predicciones precisas. En nuestro caso, lo usamos para predecir quién es \textbf{riesgo de no pagar} un crédito.

\subsection{La Idea Principal}

Imagine que usted es un experto en análisis de crédito con 30 años de experiencia. Ahora imagine que su colega le pregunta:

\begin{center}
\textit{``¿Esta persona pagará su crédito o no?''}
\end{center}

Usted miraría varios factores:
\begin{itemize}
    \item \textbf{Edad:} ¿Es joven o adulto? (experiencia de vida)
    \item \textbf{Ingreso:} ¿Tiene estabilidad financiera?
    \item \textbf{Monto del préstamo:} ¿Es proporcional a su ingreso?
    \item \textbf{Tasa de interés:} ¿Puede pagar los intereses?
    \item \textbf{Historial de crédito:} ¿Ha pagado antes?
    \item \textbf{Historial de impagos:} ¿Ha fallado en pagos anteriores?
\end{itemize}

XGBoost hace exactamente esto: \textbf{combina múltiples factores para tomar una decisión}.

---

% ============================================================================
\section{Conceptos Fundamentales}
% ============================================================================

\subsection{Árbol de Decisión: El Bloque Básico}

XGBoost está construido sobre \textbf{árboles de decisión}. Un árbol es simple: es una serie de preguntas que conducen a una respuesta.

\textbf{Ejemplo visual de un árbol simple:}

[Estructura: Ingreso > $40k? --> Edad > 30? --> Decisiones finales]

\subsection{El Truco de XGBoost: Muchos Árboles Débiles}

En lugar de crear UN árbol muy complejo (que podría cometer errores), XGBoost crea \textbf{200 árboles pequeños} que \textbf{votan juntos}.

\begin{center}
\textbf{Analogía:}
\end{center}

¿A quién le cree más?
\begin{itemize}
    \item ¿A UN experto que está 70\% seguro?
    \item ¿O a 10 expertos que CADA UNO está 60\% seguro?
\end{itemize}

\textbf{Respuesta:} A los 10 expertos. XGBoost usa esta idea: \textbf{200 pequeños árboles votan y deciden juntos}.

---

% ============================================================================
\section{Flujo de Datos: Ejemplo Paso a Paso}
% ============================================================================

\subsection{Datos de Entrada}

Tomemos 5 clientes reales de nuestro dataset:

\begin{table}[H]
\centering
\begin{tabular}{|c|c|c|c|c|c|c|c|}
\hline
\textbf{ID} & \textbf{Edad} & \textbf{Ingreso} & \textbf{Monto} & \textbf{Tasa} & \textbf{Emp. Años} & \textbf{Hist. Crédito} & \textbf{Impago Anterior} \\
\hline
1 & 35 & 50,000 & 10,000 & 8.5\% & 5 & 10 & NO \\
\hline
2 & 28 & 35,000 & 15,000 & 12.5\% & 1 & 3 & SÍ \\
\hline
3 & 45 & 25,000 & 20,000 & 15\% & 0.5 & 1 & SÍ \\
\hline
4 & 55 & 80,000 & 8,000 & 6\% & 20 & 25 & NO \\
\hline
5 & 32 & 42,000 & 12,000 & 9\% & 3 & 7 & NO \\
\hline
\end{tabular}
\caption{Ejemplo de clientes para análisis}
\end{table}

\textbf{Objetivo:} Para cada cliente, predecir si será \textbf{Riesgo (1)} o \textbf{No Riesgo (0)}.

\subsection{Paso 1: Normalización (Poner todo en la misma escala)}

Los datos vienen en diferentes escalas:
\begin{itemize}
    \item Edad: 20-70 años
    \item Ingreso: \$10,000 - \$200,000
    \item Tasa de interés: 3\% - 20\%
\end{itemize}

XGBoost necesita que TODO esté en la misma escala. Usamos \textbf{StandardScaler}:

\[
x_{\text{normalizado}} = \frac{x - \mu}{\sigma}
\]

Donde:
\begin{itemize}
    \item $x$ = valor original
    \item $\mu$ = promedio de todos los valores
    \item $\sigma$ = desviación estándar
\end{itemize}

\textbf{Ejemplo con Edad del Cliente 1:}

Supongamos que todas las edades tienen promedio $\mu = 40$ y desviación $\sigma = 10$:

\[
\text{Edad normalizada} = \frac{35 - 40}{10} = \frac{-5}{10} = -0.5
\]

Después de normalizar, \textbf{todos los datos están entre -3 y 3 aproximadamente}.

\subsection{Paso 2: Ingeniería de Características (Crear nuevas columnas inteligentes)}

XGBoost trabaja mejor si le damos \textbf{características inteligentes}. No solo los datos crudos, sino combinaciones que tengan sentido:

\begin{table}[H]
\centering
\begin{tabular}{|l|l|l|}
\hline
\textbf{Nombre} & \textbf{Fórmula} & \textbf{Interpretación} \\
\hline
Deuda/Ingreso & $\frac{\text{Monto}}{\text{Ingreso}}$ & ¿Cuánto pide vs. cuánto gana? \\
\hline
Riesgo Interés & Tasa$^2$ & Penaliza tasas altas más \\
\hline
Interacción Préstamo-Ingreso & $\frac{\text{Monto}}{\text{Ingreso} + 1}$ & Proporción normalizada \\
\hline
Estabilidad Empleo & $\log(\text{Años} + 1)$ & Experiencia laboral \\
\hline
Proporción Historial & $\frac{\text{Hist. Crédito}}{\text{Edad} + 1}$ & Historial relativo a edad \\
\hline
\end{tabular}
\caption{Nuevas características creadas para mejorar predicción}
\end{table}

\textbf{Ejemplo: Cliente 1}

\begin{align*}
\text{Deuda/Ingreso} &= \frac{10,000}{50,000} = 0.20 \quad (\text{20\% de su ingreso})\\
\text{Riesgo Interés} &= 8.5^2 = 72.25\\
\text{Est. Empleo} &= \log(5 + 1) = \log(6) \approx 1.79
\end{align*}

\subsection{Paso 3: Entrenamiento (Ajustar los Árboles)}

XGBoost crea 200 árboles pequeños. Cada árbol aprende a partir de los \textbf{errores del árbol anterior}.

\textbf{Concepto clave: Boosting (Impulsión)}

Imagine que está enseñando a un estudiante:

\begin{enumerate}
    \item Primer árbol: Intenta predecir. Comete errores.
    \item Segundo árbol: Se enfoca en los clientes que el primero predijo mal.
    \item Tercer árbol: Se enfoca en los errores restantes.
    \item ... (esto se repite 200 veces)
\end{enumerate}

Cada árbol nuevo \textbf{``corrige''} a los anteriores.

\subsection{Paso 4: Predicción para el Cliente 1}

Una vez entrenado, el modelo hace predicciones:

\begin{enumerate}
    \item \textbf{Árbol 1} vota: ``Este cliente tiene 60\% probabilidad de NO impago''
    \item \textbf{Árbol 2} vota: ``Basándome en su estabilidad laboral, 40\% probabilidad de NO impago''
    \item \textbf{Árbol 3} vota: ``Basándome en su historial, 35\% probabilidad''
    \item ... (200 votos totales)
\end{enumerate}

\textbf{Predicción Final (Promedio de todos los votos):}

\[
\text{Probabilidad de Impago} = \frac{60\% + 40\% + 35\% + \ldots}{200} \approx 15\%
\]

\textbf{Decisión:}
\begin{itemize}
    \item Si probabilidad $> 11.6\%$ (umbral óptimo) → \textbf{RIESGO}
    \item Si probabilidad $\leq 11.6\%$ → \textbf{NO RIESGO}
\end{itemize}

Para el Cliente 1: $15\% > 11.6\%$ → \textbf{Predicción: RIESGO}

---

% ============================================================================
\section{Ejemplo Numérico Completo}
% ============================================================================

\subsection{Cliente 3: El Caso de Alto Riesgo}

\begin{table}[H]
\centering
\begin{tabular}{|l|r|}
\hline
\textbf{Característica} & \textbf{Valor} \\
\hline
Edad & 45 años \\
Ingreso & \$25,000 \\
Monto del Préstamo & \$20,000 \\
Tasa de Interés & 15\% \\
Años de Empleo & 0.5 \\
Años de Historial & 1 \\
Impago Anterior & SÍ \\
\hline
\end{tabular}
\caption{Datos del Cliente 3}
\end{table}

\subsubsection{Cálculos de Características}

\begin{align*}
\text{Deuda/Ingreso} &= \frac{20,000}{25,000} = 0.80 \quad (\text{80\% de su ingreso - RIESGO ALTO!})\\
\text{Riesgo Interés} &= 15^2 = 225 \quad (\text{Tasa muy alta})\\
\text{Est. Empleo} &= \log(0.5 + 1) = \log(1.5) \approx 0.41 \quad (\text{Muy nuevo en el trabajo})\\
\text{Proporción Historial} &= \frac{1}{45 + 1} = 0.022 \quad (\text{Muy poco historial})
\end{align*}

\subsubsection{Votación de los Árboles}

\begin{table}[H]
\centering
\begin{tabular}{|c|c|l|}
\hline
\textbf{Árbol} & \textbf{Voto} & \textbf{Razón} \\
\hline
1 & 85\% & Deuda/Ingreso muy alta (0.80) \\
2 & 78\% & Tasa de interés muy alta (15\%) \\
3 & 72\% & Poco tiempo en el empleo (0.5 años) \\
4 & 88\% & Impago anterior + historial corto \\
5 & 81\% & Edad 45 con muy poco historial \\
\vdots & \vdots & \vdots \\
200 & 76\% & Múltiples factores de riesgo \\
\hline
\end{tabular}
\caption{Ejemplos de votos de algunos árboles}
\end{table}

\subsubsection{Predicción Final}

\[
\text{Probabilidad Promedio} = \frac{85\% + 78\% + 72\% + 88\% + 81\% + \ldots + 76\%}{200} \approx 80\%
\]

\textbf{Comparación con umbral:}
\[
80\% > 11.6\% \Rightarrow \textbf{PREDICCIÓN: RIESGO ALTO}
\]

\textbf{Explicación para el banco:}
\begin{quote}
Este cliente presenta múltiples señales de alerta:
\begin{itemize}
    \item Pide un monto (20k) equivalente al 80\% de su ingreso anual
    \item Tasa de interés muy alta (15\%), lo que sugiere que es cliente de alto riesgo
    \item Acaba de empezar el trabajo (0.5 años) - inestable
    \item Tiene historial de impagos
    \item \textbf{Recomendación: RECHAZAR crédito o cobrar tasa más alta}
\end{itemize}
\end{quote}

---

% ============================================================================
\section{Métricas de Desempeño}
% ============================================================================

\subsection{¿Qué Tan Bueno es Nuestro Modelo XGBoost?}

\begin{table}[H]
\centering
\begin{tabular}{|l|c|l|}
\hline
\textbf{Métrica} & \textbf{Valor} & \textbf{Significado} \\
\hline
Precisión & 82.4\% & De 100 que predice que IMPAGARÁ, 82 realmente impagan \\
\hline
Recall & 65.2\% & De todos los que realmente IMPAGAN, detectamos 65 de cada 100 \\
\hline
F1-Score & 0.73 & Balance entre Precisión y Recall (escala 0-1) \\
\hline
AUC-ROC & 0.8983 & Probabilidad de que el modelo diferencie riesgos vs no-riesgos: 89.83\% \\
\hline
\end{tabular}
\caption{Desempeño del Modelo XGBoost en conjunto de prueba}
\end{table}

\textbf{Interpretación:}
\begin{itemize}
    \item El modelo es \textbf{bastante confiable}
    \item De cada 100 clientes de riesgo real, detecta ~65
    \item La tasa de falsos positivos es aceptable para el banco
\end{itemize}

---

% ============================================================================
\section{Umbral Óptimo: El Punto de Corte}
% ============================================================================

\subsection{¿Por Qué 11.6\%?}

XGBoost da una probabilidad (0-100\%). Pero ¿cuándo decimos ``RIESGO''? ¿En 50\%? ¿En 20\%?

\textbf{Análisis de Costo-Beneficio:}

\begin{table}[H]
\centering
\begin{tabular}{|l|r|r|}
\hline
\textbf{Escenario} & \textbf{Costo} & \textbf{Valor Recuperado} \\
\hline
Rechazo Correcto (evitar impago) & \$0 & \$30,000 \\
Aceptación Correcta (cliente paga) & \$0 & \$5,000 \\
Rechazo Incorrecto (rechazamos buen cliente) & \$800 & \$0 \\
Aceptación Incorrecta (cliente impaga) & \$30,000 & \$0 \\
\hline
\end{tabular}
\caption{Matriz de Costo del Banco}
\end{table}

Con estos costos, el banco determinó que el \textbf{umbral óptimo es 11.6\%}:

\begin{itemize}
    \item Si probabilidad $\geq 11.6\%$ → Rechazar (costo de rechazo incorrecto: \$800)
    \item Si probabilidad $< 11.6\%$ → Aceptar (riesgo de impago: \$30,000)
\end{itemize}

---

% ============================================================================
\section{Ventajas del Modelo XGBoost}
% ============================================================================

\subsection{¿Por Qué Usamos XGBoost en Lugar de Otros Métodos?}

\begin{enumerate}
    \item \textbf{Precisión:} 89.83\% AUC-ROC (muy preciso)
    
    \item \textbf{Maneja Datos Complejos:} Puede capturar relaciones no-lineales entre variables
    
    \item \textbf{Rápido:} Predicciones en milisegundos
    
    \item \textbf{Robusto:} Funciona bien incluso con datos imperfectos
    
    \item \textbf{Escalable:} Puede procesar millones de clientes
    
    \item \textbf{Interpretable:} Podemos explicar por qué rechaza o acepta
\end{enumerate}

---

% ============================================================================
\section{Resumen del Flujo Completo}

[Flujo: Datos Cliente --> Normalizar --> Ingeniería Características --> 200 Árboles --> Promedio Votos --> Umbral 11.6\% --> Predicción]

---

% ============================================================================
\section{Conclusión}
% ============================================================================

XGBoost es un método poderoso que:

\begin{enumerate}
    \item Toma \textbf{múltiples factores} de un cliente
    \item Los \textbf{transforma inteligentemente} en características
    \item Los \textbf{pasa por 200 árboles} que votan
    \item \textbf{Combina los votos} para una predicción
    \item \textbf{Compara con un umbral óptimo} para tomar la decisión final
\end{enumerate}

\textbf{Resultado final:} Un sistema \textbf{confiable al 89.83\%} para identificar clientes de riesgo, permitiendo al banco:
\begin{itemize}
    \item Reducir pérdidas por impagos
    \item Maximizar utilidades en colocaciones seguras
    \item Tomar decisiones rápidas y consistentes
\end{itemize}

---

\appendix

\section{Glosario de Términos}

\begin{description}
    \item[Árbol de Decisión:] Modelo que hace predicciones preguntando preguntas sucesivas
    \item[Boosting:] Técnica donde múltiples modelos débiles se combinan para crear uno fuerte
    \item[AUC-ROC:] Métrica que mide qué tan bien el modelo diferencia entre riesgos y no-riesgos
    \item[Normalización:] Poner todos los datos en la misma escala
    \item[Ingeniería de Características:] Crear nuevas variables que tengan sentido para el problema
    \item[Umbral:] Punto de corte para tomar una decisión
\end{description}

\end{document}
